\documentclass[11pt,a4paper]{article}
\usepackage[utf8]{inputenc}
\usepackage[T1]{fontenc}
\usepackage[french]{babel}
\usepackage{geometry}
\usepackage{xcolor}
\usepackage{hyperref}
\usepackage{titlesec}

\geometry{margin=1in}
\definecolor{titleblue}{RGB}{0,51,102}

\hypersetup{
    colorlinks=true,
    linkcolor=titleblue,
    urlcolor=titleblue,
    pdftitle={Lettre de Motivation - Ismail AISSAOUI},
}

\begin{document}

\noindent \textbf{Ismail AISSAOUI} \\
Ingénieur d'État en Intelligence Artificielle (ESI-SBA) \\
Email: ismail.aissaoui.pro@gmail.com \\
Téléphone: +213 660 70 77 96 \\
LinkedIn: https://ismail-aissaoui.vercel.app

\bigskip

\begin{flushright}
    Au Comité de Sélection du Doctorat \\
    \textbf{Programme EDT - Inria Rennes} \\
    À l'attention de : Prof. Benoit Caillaud \\
    \medskip
    1er mai 2026
\end{flushright}

\bigskip

\noindent \textbf{Objet : Candidature au doctorat : Méthodes structurelles pour l'ingénierie de Jumeaux Numériques mixtes modèles/données (Rang 2 - PC1)}

\bigskip

Monsieur le Professeur,

Ingénieur d'État en Intelligence Artificielle diplômé de l'ESI-SBA, je vous adresse avec enthousiasme ma candidature pour le poste de doctorat intitulé « Méthodes structurelles pour l'ingénierie de Jumeaux Numériques mixtes modèles/données » au sein du programme national Engineering Digital Twin (EDT). Cette candidature est motivée par la forte adéquation entre mon expertise en diagnostics basés sur les GNN et vos recherches sur l'analyse structurelle des systèmes DAE.

Actuellement en fin de cycle ingénieur chez Sonatrach, je travaille sur le développement de **modèles de substitution informés par la physique** pour la surveillance de turbines à gaz. Mes recherches portent sur l'intégration de lois physiques dans des architectures de deep learning pour combler le « fossé de réalité » (reality gap). J'ai acquis une solide expérience dans l'utilisation des **Réseaux de Neurones sur Graphes (GNN)** pour modéliser les dépendances structurelles au sein de systèmes complexes, une compétence que je souhaite appliquer à l'analyse structurelle des systèmes d'équations algébro-différentielles (DAE).

Le projet que vous proposez m'intéresse particulièrement car il combine la rigueur de l'analyse structurelle (détection de sous-systèmes minimalement surdéterminés - MSO) avec la flexibilité des approches basées sur les données. Je suis convaincu que mon expérience en implémentation PyTorch et ma compréhension des problématiques de diagnostic industriel me permettront de contribuer de manière significative aux travaux de l'équipe-projet à l'Inria Rennes.

Ma maîtrise des outils de simulation et mon autonomie me permettront de m'intégrer rapidement au consortium EDT et de contribuer au développement de la **plateforme Artemis**.

Je vous remercie de l'attention que vous porterez à ma candidature. Je reste à votre entière disposition pour un entretien afin de discuter de mon projet de recherche et de ma motivation pour rejoindre votre équipe.

Dans l'attente de votre réponse, je vous prie d'agréer, Monsieur le Professeur, l'expression de mes salutations distinguées.

\bigskip

\noindent Ismail AISSAOUI

\end{document}
